\documentclass[11pt]{article}
\usepackage[T1]{fontenc}
\usepackage[margin=1in]{geometry}
\usepackage{amsmath,amssymb,amsthm,mathtools}
\usepackage[colorlinks=true,linkcolor=blue,urlcolor=blue,citecolor=blue]{hyperref}
\usepackage{microtype}
\newtheorem{theorem}{Theorem}[section]
\newtheorem{lemma}[theorem]{Lemma}
\newtheorem{proposition}[theorem]{Proposition}
\newtheorem{corollary}[theorem]{Corollary}
\theoremstyle{definition}
\newtheorem{definition}[theorem]{Definition}
\theoremstyle{remark}
\newtheorem{remark}[theorem]{Remark}
\newcommand{\one}{\mathbf 1}
\newcommand{\E}{\mathbb E}
\newcommand{\dd}{\,\mathrm d}
\newcommand{\Tr}{\operatorname{Tr}}
\newcommand{\SU}{\mathrm{SU}(2)}
\newcommand{\norm}[1]{\left\lVert#1\right\rVert}
\newcommand{\osc}{\operatorname{osc}}
\setlength{\emergencystretch}{3em}
\hypersetup{pdftitle={Weighted boundary costs and accepted local repairs},pdfauthor={Research note},pdfsubject={A scoped finite lattice comparison, not a Yang--Mills or RG construction}}
\title{Weighted boundary costs and accepted local repairs}
\author{Research note}
\date{October 9, 2026}
\begin{document}
\maketitle
\begin{abstract}
We separate two requirements for weighted rough-region estimates:
control of a weight's change during an edit and control of its bias on
frozen boundary labels. The former does not imply the latter, even for
positive, local, gauge-invariant weights: an exterior-only tilt can make
a rare frustrated boundary likely while remaining unchanged by the edit.
We prove a one-region comparison under an explicit local marginal-density
condition, and verify a different sufficient structural class without
assuming that density condition: bounded reflection-symmetric interactions
on elementary four-cells. In this class, weighted boundary-energy moments
and local edit costs have constants proportional to the affected region,
not the full volume. An accepted low-minimum repair has an output
submeasure dominated by a small multiple of the same weighted measure;
this permits composition of specified successful repairs. Exceptional
retained labels do not admit that multiplication rule. All conclusions
are for positive finite-lattice base measures. We do not establish a
mandatory covering expansion, RG closure, cutoff comparison, or a
continuum Yang--Mills theory. No novelty claim is made. The third-pass
input has been adversarially reread and its geometric application made
explicit; this is not independent expert or agent certification.
\end{abstract}
\tableofcontents

\section{Definitions and the third-pass input}
Use $\Lambda_L=(\mathbb Z/L\mathbb Z)^4$, $L=2^m$, $m\ge2$, with
$V=L^4$ sites, $4V$ links, and $6V$ plaquettes. Orient each link positively
and assign $U_e\in\SU$, with the inverse on its reverse. The reference
measure is product probability Haar. For $H\ge1$ put
\[
 s(g)=1-\tfrac12\Tr g=\tfrac12\norm{\one-g}_{\rm op}^2,\quad
 E=\sum_qs(U_q),\quad
 \mu_H(\dd U)=Z_H^{-1}e^{-HE(U)}\dd U.
\]
For a fixed nonempty editable link set $J$, let $n=|J|$, let $P(J)$
contain every plaquette touching $J$, and set
\[
 E_J=\sum_{q\in P(J)}s(U_q),\qquad
 m_J(z)=\min_{x\in\SU^J}E_J(x,z),\quad z=U_{J^c}.
\]
Let $F_J$ be the frozen links appearing on these plaquettes and
$\mathcal F_J=\sigma(U_e:e\in F_J)$. Then $m_J$ is continuous and
$\mathcal F_J$-measurable, with $|F_J|\le18n$. This local sigma algebra
is smaller than that of the complete frozen exterior.

Let $N_p(J)$ count plaquettes in $P(J)$ with operator curvature greater
than $p/\sqrt H$, where $p>0$. Define
\begin{equation}
 \ell_H=\log24+\tfrac32\log H,\quad
 a_H=\tfrac32+\log24+\tfrac32\log(8H),\quad
 b_H=\tfrac12a_H+\frac{\log2}{64}.
 \label{eq:constants}
\end{equation}
The third-pass note \cite{third} proves
\begin{equation}
 \E_{\mu_H}e^{HE_J/128}\le e^{b_Hn},\qquad
 \mu_H(m_J>v)\le\min\{1,e^{-Hv/128+b_Hn}\}.
 \label{eq:raw}
\end{equation}
It also constructs a Borel gauge-covariant kernel $K_J$ which freezes
$J^c$ and samples probability Haar restricted to
\[
 B_J(z)=\{y:E_J(y,z)\le m_J(z)+12n/H\}.
\]
Its edited-link density is at most $M_J=e^{\ell_Hn}$. The minimizer-cap
argument ensures the target has positive Haar volume at least $M_J^{-1}$;
no measurable choice of a minimizer is needed. This is an existence
construction, not an efficient algorithm for sampling or optimization.

For a positive integer $k$ write
\begin{equation}
 \begin{split}
 R_{J,k}&=\{N_p(J)\ge k\},\qquad u=kp^2/(4H),\\
 A_{J,k}&=R_{J,k}\cap\{m_J\le u\},\qquad
 B_{J,k}=\{m_J>u\}.
 \end{split}
 \label{eq:events}
\end{equation}
Thus $R_{J,k}\subset A_{J,k}\cup B_{J,k}$. If $kp^2\ge96n$, every
output from $A_{J,k}$ obeys
\begin{equation}
 N_p(J;y)\le3k/4,\qquad H(E(x)-E(y))\ge kp^2/8.
 \label{eq:gain}
\end{equation}
This reduces the total affected rough count, not necessarily each old
rough plaquette; individual newly rough plaquettes are allowed.

\section{Adversarial audit of the chessboard application}
\label{sec:audit}
The analytic input behind \eqref{eq:raw} is the standard event chessboard
inequality \cite[\S2.2, Theorem 2.4]{bk}. The closed-block convention and
the fact that only the designated block-boundary reflections are needed
are also explicit in \cite[\S5.3, Theorem 5.8]{notes}. We retain the
power-of-two restriction, and use no assertion of independent blocks.

Here is a precise geometric encoding for the link-variable application.
Put the variable for $(x,x+e_\mu)$ at $2x+e_\mu$ on the refined torus
$(\mathbb Z/(2L)\mathbb Z)^4$. These are exactly the points with one odd
coordinate and the other three even. Fill the other refined sites with
deterministic dummy variables. An original reflection
$x_i\mapsto2a-x_i$ becomes $v_i\mapsto4a-v_i$. It preserves this parity
pattern, swaps the two closed-half edge algebras, and inverts the group
variable if its odd coordinate is $i$. Tangential fixed-plane variables
are pointwise fixed. Reflections in different coordinate directions
commute, including their orientation inversions. No commutation of
parallel reflections at different positions is asserted.

For the Wilson action no plaquette meets both strict halves. With $z$
the fixed-plane links and $x_+,x_-$ the remaining variables, the energy
splits as $E_P(z)+E_+(x_+,z)+E_-(x_-,z)$ with the two half-energies
reflections of each other. Thus for a bounded real positive-half function
$f$, its reflection $\vartheta f$ satisfies
\begin{equation}
 \E_{\mu_H}[f\vartheta f]
 =Z_H^{-1}\int e^{-HE_P(z)}
   \left(\int e^{-HE_+(x_+,z)}f(x_+,z)\dd x_+\right)^2\dd z\ge0.
 \label{eq:rp}
\end{equation}
Original side-two blocks become closed side-four refined blocks; the
factor torus still has $M=V/16$ blocks. Their shared boundary variables
are precisely the variables kept as $z$, not independent replicas. This
checks the support and reflection-algebra assumptions used in the
repeated Cauchy--Schwarz proof of the event chessboard estimate.

A prototype link has local starting coordinate 0 or 1 along its direction
and 1 in the other three coordinates. All six star plaquettes have
relative interiors inside its side-two block. Reflected stars in
different blocks therefore have disjoint plaquette sets. The 64 classes
are direction plus four starting-site parities. For each class the grid
can be shifted so its links occupy distinct blocks; an odd longitudinal
block changes the pulled-back prototype from 0 to 1. Both prototypes
satisfy the same disjointness bound. This verifies the hypotheses for
unequal thresholds at different stars.

Finally, a product cap of radius $(8H)^{-1/2}$ gives
$Z_H\ge e^{-4Va_H}$. A disseminated event requiring star energy at least
$v$ forces total energy at least $Mv$, hence has chessboard weight at
most $\min\{1,e^{-Hv+64a_H}\}$. Within a class, joint tails factorize
\emph{as an inequality}. Layer-cake integration and H\"older over the
64 classes give \eqref{eq:raw}. The four occurrences of each plaquette
in the uncolored sum of all stars are not discarded as if disjoint.

No mathematical error requiring a change to the third-pass proof was
found in this bounded reread. Its restricted volume class, raw-marginal
scope, overlap losses, and distinction between output-weighted integrals
and probabilities remain necessary. This is a same-agent audit, not an
independent certification. The accompanying exact finite-model and
geometric tests are additional diagnostics, not proofs of the analytic
chessboard theorem.

\section{The missing marginal condition and a counterexample}
Let $Q>0$ be finite almost everywhere and satisfy
$0<\E_{\mu_H}Q<\infty$. Define
\[
 \nu_Q(\dd U)=\frac{Q(U)}{\E_{\mu_H}Q}\mu_H(\dd U),\qquad
 r_J=\frac{\E_{\mu_H}[Q\mid\mathcal F_J]}{\E_{\mu_H}Q}.
\]
Then $r_J$ is exactly the Radon--Nikodym density of the weighted
$F_J$ marginal relative to the raw one. For any $\mathcal F_J$-measurable
event $B$,
\begin{equation}
 \nu_Q(B)=\E_{\mu_H}[r_J\one_B].
 \label{eq:marginal}
\end{equation}
The normalization is the \emph{full} expectation of $Q$. A conditional or
restricted normalizer cannot be substituted without changing this law.

\begin{proposition}[Edit stability alone gives no exceptional-label cost]
\label{prop:counterexample}
There are strictly positive local gauge-invariant weights, constant on
every $J$-edit, which assign any prescribed probability $s\in(0,1)$ to
a given high-minimum boundary event of raw probability $t\in(0,1)$.
\end{proposition}
\begin{proof}
For one link take $B=\{m_J>v\}$ with $0<v<6$. The balanced frustrated
star has $m_J=6$, identity staple data have $m_J=0$, and both have open
neighborhoods on the corresponding side of the threshold. The finite
Wilson law has everywhere positive continuous density, so $0<t=\mu_H(B)<1$.
The event depends only on the 18 frozen staple links and is gauge invariant.
Set
\[
 Q=q\one_B+\one_{B^c},\qquad q=\frac{s(1-t)}{(1-s)t}.
\]
This weight is positive, finite and unchanged by any edited central-link
value, but $\nu_Q(B)=qt/(qt+1-t)=s$. Its normalized local marginal density
is $s/t$ on $B$ and $(1-s)/(1-t)$ on $B^c$. Hence, for $r>1$,
\begin{equation}
 \norm{r_J}_{L^r(\mu_H)}^r
 =s^r t^{1-r}+(1-s)^r(1-t)^{1-r}.
 \label{eq:sharpness}
\end{equation}
The missing marginal cost becomes large as $t$ decreases at fixed $s$,
although the edit ratio stays exactly one.
\end{proof}
The phenomenon is not an artifact of discontinuous weights. Choose a
continuous function $f(m)$ that vanishes for $m\le v$, equals one on
$m\ge v'>v$ with $v'<6$, and lies in $[0,1]$. Then $Q_T=e^{Tf(m_J)}$
is continuous, positive, gauge invariant and edit-invariant, while
$\nu_{Q_T}(B)\to1$ as $T\to\infty$: a positive-measure subset of $B$
has weight $e^T$, while $B^c$ has weight one. No uniform amplitude bound
in $T$ or $H$ is claimed.

\begin{lemma}[Local marginal transfer]\label{lem:holder}
For $1<r\le\infty$, let $\alpha=1-1/r$ (with $1/\infty=0$). If
$\norm{r_J}_{L^r(\mu_H)}\le e^{c n}$, then
\begin{equation}
 \nu_Q(B_{J,k})\le
 \min\{1,e^{cn}\mu_H(B_{J,k})^\alpha\}
 \le\min\{1,\exp[-\alpha kp^2/512+(c+\alpha b_H)n]\}.
 \label{eq:holder}
\end{equation}
\end{lemma}
\begin{proof}
Apply H\"older to \eqref{eq:marginal}, and then \eqref{eq:raw} at
$v=kp^2/(4H)$. Using the bound without its truncation by one is valid;
the final probability can again be truncated by one.
\end{proof}
This is an explicit sufficient hypothesis, not a proof that an arbitrary
RG-generated weight satisfies it. It concerns a local projection of $Q$,
not its whole-volume $L^r$ norm. Equation~\eqref{eq:sharpness} explains
why some quantitative control of boundary bias is indispensable.

A directly verified elementary class is $Q=e^W$ with
$\sup W-\inf W\le c n$: then $r_J\le e^{cn}$ and every edit ratio is at
most $e^{cn}$. For example a perturbation supported on at most $s n$
links, with total log-oscillation at most $c_0s n$, has this property.
\emph{Local support alone} is not the bound, and adding an arbitrary
exterior factor destroys the argument. An extensive bounded interaction
usually has total oscillation of order $V$; the next structural class
avoids importing that whole-volume loss into a local estimate.

\section{An output domination lemma for positive weights}
For a source event $A$, write
\[
 T_{J,A}(x,\dd y)=\one_A(x)K_J(x,\dd y).
\]
This is an accepted-transition subprobability kernel, not the Markov
kernel obtained by staying put on rejected sources.

\begin{lemma}[Weighted accepted output]\label{lem:domination}
Suppose on source/kernel-output pairs from $A$,
\[
 H(E(x)-E(y))\ge g,\qquad Q(x)\le e^{d n}Q(y).
\]
Then as measures,
\begin{equation}
 \nu_Q T_{J,A}\le C_J\nu_Q,\qquad
 C_J=\exp[-g+(\ell_H+d)n].
 \label{eq:domination}
\end{equation}
In particular $\nu_Q(A)\le\min\{1,C_J\}$, and for every measurable $f$
with $\E_{\nu_Q}|f|<\infty$,
\begin{equation}
 \frac1{Z_Q}\int_A e^{-HE(x)}Q(x)(K_J|f|)(x)\dd x
 \le C_J\E_{\nu_Q}|f|,
 \quad Z_Q=\int e^{-HE}Q.
 \label{eq:observable}
\end{equation}
\end{lemma}
\begin{proof}
Write $x=(\xi,z)$, $y=(\eta,z)$ and let $k_z(\xi,\eta)\le M_J$ be
the edited-output density. On the indicated pairs,
$e^{-HE(\xi,z)}Q(\xi,z)\le e^{-g+dn}e^{-HE(\eta,z)}Q(\eta,z)$.
For a nonnegative output test function $f$, insert this inequality, bound
$k_z$ by $M_J$, and integrate the original $\xi$ over probability Haar.
Its mass is at most one. Tonelli leaves at most
$e^{-g+dn}M_J\int f\dd\nu_Q$, proving measure domination. Since $K_J$
is a probability kernel, the total mass on the left is $\nu_Q(A)$.
Apply the result to $|f|$ for \eqref{eq:observable}.
\end{proof}
The positive base weight remains at the actual output after the
comparison. Signed or complex \emph{test functions} may be estimated by
their absolute values; signed/complex base densities are not allowed.
For contrast the earlier literal integral
$Z_Q^{-1}\int_Ae^{-HE(x)}(K_JQ)(x)\dd x$ has a bound without $e^{dn}$,
but is not in general $\nu_Q(A)$. These expressions are not identified.

\begin{theorem}[One-region comparison with declared weight hypotheses]
\label{thm:general}
Assume $kp^2\ge96n$, the edit-ratio bound of
Lemma~\ref{lem:domination} on $A_{J,k}$, and the local marginal bound of
Lemma~\ref{lem:holder}. Then
\begin{equation}
 \nu_Q(R_{J,k})\le\min\left\{1,
 e^{-kp^2/8+(\ell_H+d)n}
 +e^{-\alpha kp^2/512+(c+\alpha b_H)n}\right\}.
 \label{eq:general}
\end{equation}
The low-minimum transition has the output domination
\eqref{eq:domination} with $g=kp^2/8$ and reduces its rough count as in
\eqref{eq:gain}.
\end{theorem}
\begin{proof}
Use $R_{J,k}\subset A_{J,k}\cup B_{J,k}$ and apply the two lemmas in the
\emph{same} measure $\nu_Q$. They control different branches, with
separate hypotheses.
\end{proof}

\section{A verified class of extensive positive perturbations}
\label{sec:cells}
For each $x\in\Lambda_L$ let $C_x=x+[0,1]^4$ be an elementary four-cell,
including its 32 links. There are $V$ such cells. Each link belongs to
exactly eight cells. We admit
\begin{equation}
 Q=e^W,\qquad W=\sum_C\phi_C,
 \label{eq:cellweight}
\end{equation}
where $\phi_C$ is a bounded real Borel function of the links in $C$,
gauge invariant, and $\osc\phi_C\le\varepsilon$. Require translation
covariance and covariance under every integer-site coordinate reflection,
including orientation reversal of links:
$\phi_{\vartheta C}(\vartheta U)=\phi_C(U)$. No smoothness or sign
restriction on the real $\phi_C$ is needed. The Boltzmann weight itself
is always positive. The amplitude $\varepsilon\ge0$ may depend on $H$;
useful small-field budgets will require it to be sufficiently small
relative to $p^2$, not necessarily an absolute small constant.

Subtract $\inf\phi_C$ from each term. These are configuration-independent
constants and do not change $\nu_Q$ or edit ratios; covariance is retained.
Thus we may assume
\begin{equation}
 0\le\phi_C\le\varepsilon,\qquad 0\le W\le\varepsilon V.
 \label{eq:range}
\end{equation}
Although the global range is extensive, the resulting local estimates
below are not multiplied by $e^{\varepsilon V}$.

\begin{lemma}[Weighted reflection and local edit cost]\label{lem:cellrp}
The measure $\nu_W=Z_W^{-1}e^{-HE+W}\dd U$ is site-reflection positive.
For any pair of configurations differing only on $J$,
\begin{equation}
 |W(x)-W(y)|\le8\varepsilon n.
 \label{eq:edit}
\end{equation}
\end{lemma}
\begin{proof}
An elementary cell has extent one in each coordinate, so its link support
cannot meet both strict halves separated by an integer-coordinate plane.
All terms in $W$ therefore split into plane and reflected half terms,
just as the Wilson plaquettes do. The positive fixed-plane factor and
the square in \eqref{eq:rp} are replaced by
\[
 e^{-HE_P(z)+W_P(z)},\qquad
 \left(\int e^{-HE_+(x_+,z)+W_+(x_+,z)}f(x_+,z)\dd x_+\right)^2.
\]
This proves positivity. Only cells containing a link of $J$ can change
under the edit, and there are at most $8n$ such cells. Sum their
oscillation bounds. No independence of cells is used.
\end{proof}

\begin{theorem}[Weighted local energy moments]\label{thm:cellmoment}
For the class \eqref{eq:cellweight}, for every fixed $J$ and $0<\theta<1$,
\begin{equation}
 \E_{\nu_W}e^{\theta HE_J/64}
 \le\exp\left[n\left(\theta(a_H+\varepsilon/4)
                      -\frac{\log(1-\theta)}{64}\right)\right].
 \label{eq:cellmoment}
\end{equation}
The same bound holds with $E_J$ replaced by $m_J$. Set
$\widehat b_H=b_H+\varepsilon/8$. Then
\begin{equation}
 \nu_W(E_J\ge v),\ \nu_W(m_J>v)
 \ \le\ \min\{1,e^{-Hv/128+\widehat b_H n}\}.
 \label{eq:celltail}
\end{equation}
These constants are uniform in the dyadic volume and all admitted
potentials with the given oscillation bound.
\end{theorem}
\begin{proof}
The centered weight satisfies $W\ge0$, so the product cap still gives
$Z_W\ge e^{-4Va_H}$. A disseminated star event forcing $E\ge Mu$ has
numerator at most $e^{-HMu+\varepsilon V}$. Its chessboard weight is
therefore at most
\begin{equation}
 q_{H,\varepsilon}(u)
 =\min\{1,\exp[-Hu+64a_H+16\varepsilon]\}.
 \label{eq:tiltedstar}
\end{equation}
Lemma~\ref{lem:cellrp} licenses the same event chessboard inequality in
$\nu_W$. The geometric disjointness and the 64 classes are unchanged.

For completeness, write $K=64a_H+16\varepsilon$ and $T_e=HE_e$.
Within a class, joint tails are bounded by products of
$q(t)=\min(1,e^{K-t})$. Expanding
$e^{\theta T_e}=1+\int_0^\infty\theta e^{\theta t}\one_{T_e>t}\dd t$
in a finite product and applying Tonelli gives
\[
 \E_{\nu_W}e^{\theta\sum_{e\in S}T_e}
 \le\left(1+\int_0^\infty\theta e^{\theta t}q(t)\dd t\right)^{|S|}
 =\left(\frac{e^{\theta K}}{1-\theta}\right)^{|S|}.
\]
Since $E_J\le\sum_{e\in J}E_e$, H\"older with 64 equal exponents proves
\eqref{eq:cellmoment}. Pointwise $m_J\le E_J$ and Markov's inequality at
$\theta=1/2$ prove \eqref{eq:celltail}.
\end{proof}
This theorem proves the needed exceptional-label tail \emph{directly in
the weighted measure}. It does not prove, and does not require, the
local likelihood-ratio hypothesis of Lemma~\ref{lem:holder} relative to
$\mu_H$. These are two distinct sufficient routes.

\begin{corollary}[Verified one-region weighted comparison]\label{cor:cellsplit}
For \eqref{eq:cellweight}, if $kp^2\ge96n$, then
\begin{equation}
 \nu_W(R_{J,k})\le\min\left\{1,
 e^{-kp^2/8+(\ell_H+8\varepsilon)n}
 +e^{-kp^2/512+\widehat b_Hn}\right\}.
 \label{eq:cellsplit}
\end{equation}
The low-minimum output submeasure is dominated by the first exponential
times $\nu_W$, for any nonnegative output test function. The exceptional
label is measured in that same $\nu_W$. If also
\begin{equation}
 kp^2\ge16n(\ell_H+8\varepsilon),\qquad
 kp^2\ge1024n\widehat b_H,
 \label{eq:budgets}
\end{equation}
then the first term is at most $e^{-kp^2/16}$ and the second at most
$e^{-kp^2/1024}$.
\end{corollary}
\begin{proof}
Combine \eqref{eq:edit} with Lemma~\ref{lem:domination} on $A_{J,k}$.
For $B_{J,k}$ apply \eqref{eq:celltail} with $v=kp^2/(4H)$.
The numerical absorptions follow directly from \eqref{eq:budgets}.
\end{proof}
As a direct probability estimate the split is weaker than the consequence
$\nu_W(R_{J,k})\le\min\{1,e^{-kp^2/256+\widehat b_Hn}\}$ of
\eqref{eq:celltail}. Its additional content is the \emph{accepted output
domination}, with a retained exceptional branch in the same measure.

\paragraph{A nontrivial explicit example.}
Take
\begin{equation}
 W_\gamma=\gamma\sum_q s(U_q)^2,\qquad\gamma\in\mathbb R.
 \label{eq:example}
\end{equation}
Each plaquette belongs to four elementary four-cells, and a cell contains
24 plaquettes. The decomposition
$\phi_C=(\gamma/4)\sum_{q\subset C}s(U_q)^2$ has
$\osc\phi_C\le24|\gamma|$, since $0\le s^2\le4$. Thus
$\widehat b_H=b_H+3|\gamma|$ is valid. Directly summing only the affected
plaquettes improves the edit constant to
\[
 |W_\gamma(x)-W_\gamma(y)|\le4|\gamma||P(J)|\le24|\gamma|n.
\]
Use this in place of $8\varepsilon n$ in \eqref{eq:cellsplit} and
\eqref{eq:budgets}. Both signs of $\gamma$ are allowed. This is not just
a change of the coefficient of $s(U_q)$ in the continuous $\SU$ model.
It is an example of the proved class, not a claim that exact RG stays
inside it.

\section{Counting and what composition actually permits}
The link-sharing-a-plaquette graph has degree at most 18. For the actual
rough-link component through a fixed root $e_0$, the same fixed-subset
counting argument as in \cite{third} now gives
\begin{equation}
 \nu_W(|C_p(e_0)|\ge n)
 \le\min\{1,18^{2(n-1)}e^{-n[p^2/1024-\widehat b_H]}\}.
 \label{eq:components}
\end{equation}
Every selected link has a rough witness, and a plaquette witnesses at
most four links; no boundary compatibility or independence is assumed.
The total weight $W$ is fixed before choosing supports. Rarity is not
re-evaluated under a separately chosen measure for each region.
For $\varepsilon_H=O(\log H)$ and $p=(\log H)^P$, $P>1/2$, the component
and one-region absorption budgets hold eventually when $k\ge n/4$.
This is only a large-$H$ arithmetic assertion in the declared model class.

\begin{proposition}[Composition of accepted repairs]\label{prop:composition}
Fix one positive base measure $\nu_Q$. Suppose nonnegative
subprobability kernels $T_i$ satisfy
$\nu_QT_i\le c_i\nu_Q$. Then
\begin{equation}
 \nu_QT_1\cdots T_r\le\left(\prod_{i=1}^r c_i\right)\nu_Q.
 \label{eq:composition}
\end{equation}
In particular this applies to the low-minimum kernels in
Corollary~\ref{cor:cellsplit} on predetermined supports, with their gates
$A_{J_i,k_i}$ evaluated on each \emph{actual intermediate state}.
No disjointness of affected plaquettes is required for this statement.
\end{proposition}
\begin{proof}
A nonnegative kernel preserves order of measures. Apply $T_2$ to
$\nu_QT_1\le c_1\nu_Q$, then use $\nu_QT_2\le c_2\nu_Q$, and iterate.
For the asserted application, every one-step energy drop is proved
against its actual source by \eqref{eq:gain}; gains are not all measured
against the common original configuration.
\end{proof}
For adaptive support choice from a fixed finite or countable family,
partition the source-selection events and use the corresponding accepted
kernels. Their sum is dominated by $\sum_Jc_J$ times the base measure.
For connected $J$ of size $n$ through a fixed root with a uniform
$k\ge n/4$ budget, there are at most $18^{2(n-1)}$ terms. This explicitly
pays the support count. A choice over arbitrary roots costs its additional
locations. A policy depending on history needs a valid bound for each
encoded history or a genuine aggregate kernel bound; the fixed-family
statement does not silently supply one.

\begin{corollary}[A counted adaptive accepted kernel]\label{cor:adaptive}
Fix a root $e_0$ and the cell-weight measure $\nu_W$. On each source choose
its current nonempty rough-link component $J=C_p(e_0)$ and set
$k_J=\lceil|J|/4\rceil$. Apply $K_J$ only when
$m_J\le k_Jp^2/(4H)$, and otherwise give this accepted kernel $T$ mass
zero. If $p^2\ge384$ and
\[
 d_H=p^2/32-\ell_H-8\varepsilon>2\log18,
\]
then
\begin{equation}
 \nu_W T\le c_H\nu_W,\qquad
 c_H=\frac{e^{-d_H}}{1-18^2e^{-d_H}}.
 \label{eq:adaptive}
\end{equation}
Thus $\nu_WT^r\le c_H^r\nu_W$. In particular $c_H<1$ if
$d_H>\log325$. This covers only paths on which all $r$ selected repairs
are accepted; the root and weight remain fixed.
\end{corollary}
\begin{proof}
The component has at least $k_J$ rough witnesses and
$k_Jp^2\ge96|J|$, so every accepted fixed-$J$ transition has domination
constant at most $e^{-d_H|J|}$. The finite collection of component-selection
events partitions the nonempty sources. Count its connected supports
through $e_0$ by $18^{2(n-1)}$ and sum the geometric series, extending the
finite size range to infinity for an upper bound. This is an aggregate
one-step kernel estimate, not a conditional Gibbs calculation. Apply
Proposition~\ref{prop:composition}. The condition for $c_H<1$ is
$325e^{-d_H}<1$.
\end{proof}
This construction recalculates the component from the actual output at
its next step. It may introduce individual new rough plaquettes and
change component geometry; the aggregate estimate already counts those
possibilities at each accepted step. It is not a mandatory-inventory
preserving repair and does not treat the retained-exception branches.

\paragraph{Unconditional sublevel output.}
For the cell class, $E(y)\le E(x)+12n/H$ holds for every source/output
pair of the ungated $K_J$. Thus the same domination proof with $g=-12n$
gives
\begin{equation}
 \nu_WK_J\le e^{(12+\ell_H+8\varepsilon)n}\nu_W.
 \label{eq:ungated}
\end{equation}
This permits specified ungated edits with an \emph{accumulating} density
cost, not cost-free reuse of the original measure.

\paragraph{Retained exceptional labels do not multiply.}
The identity branch $S_B(x,\dd y)=\one_B(x)\delta_x(\dd y)$ obeys
$\nu_QS_B=\one_B\nu_Q$, which is not dominated by
$\nu_Q(B)\nu_Q$ if $0<\nu_Q(B)<1$. Applying it twice still has mass
$\nu_Q(B)$, not $\nu_Q(B)^2$. A small marginal probability is not the
small \emph{operator} domination needed in
Proposition~\ref{prop:composition}. Joint exceptional events need a
joint-energy/overlap estimate or additional decorrelation; no product of
one-label tails is asserted. Even disjoint link edits can share affected
plaquettes, and a history can revisit the same boundary defect.

For fixed sets $J_i$, one available joint input remains
\[
 \E_{\nu_W}\exp\left[\frac H{128m_*}\sum_i m_{J_i}\right]
 \le e^{\widehat b_H|\bigcup_iJ_i|},\qquad
 m_* =\max_q|\{i:q\in P(J_i)\}|,
\]
by $\sum_i m_{J_i}\le m_*E_{\cup_iJ_i}$ and
Theorem~\ref{thm:cellmoment}. It concerns labels evaluated on \emph{one}
configuration, not a sequence of changed exteriors.

\section{Limits and the next precise obligation}
\paragraph{Why the hypotheses cannot be discarded.}
Proposition~\ref{prop:counterexample} disproves an exceptional-label
bound from local edit stability alone. Reflection symmetry and positivity
of a weight alone also do not imply reflection positivity: two symmetric
spins $\sigma_\pm\in\{-1,1\}$ with weight
$e^{-\tau\sigma_+\sigma_-}$, $\tau>0$, have
$\E[\sigma_+\vartheta\sigma_+]=-\tanh\tau<0$. The interaction crosses
both strict halves. This simple example is not a counterexample to the
cell-local theorem; it explains why a longer interaction needs its own
positivity proof.

The weight $Q=e^{HE}$ cancels Wilson suppression. It even has a cell
representation with $\varepsilon=12H$, by distributing plaquette energies
over their four cells. Our resulting constants are then too large to
imply small rough probabilities: $\widehat b_H\ge3H/2$, while
$p^2<4H$ in the nonempty rough regime. Thus the admitted-class theorem
does not claim a useful bound at arbitrary perturbation amplitudes.

The frustrated star, padded isolated defect, flat winding holonomy, and
the overlapping-edit example from the preceding notes remain valid.
Neither rarity nor a reduction in rough count implies an all-small
boundary-compatible filling or preservation of a prescribed inventory.
The probability statements remain restricted to periodic dyadic cubes,
not arbitrary frozen boundary conditions or all Gibbs states.

\paragraph{What is now discharged.}
For the explicitly admitted positive cell-interaction class, both branches
of a one-region comparison are controlled in the \emph{same} weighted
measure with region-proportional costs. The accepted low-minimum output
has enough measure domination to compose specified successful repairs.
For general positive weights the exact missing local marginal-density
hypothesis is isolated and shown not to follow from edit stability.

\paragraph{What remains.}
The spin-model endpoint comparison inspected in \cite[\S\textup{Integrated
comparison of endpoint densities}, equations labeled
\texttt{eq:orig-16} and \texttt{eq:orig-22}]{o3} uses a mandatory covering
sum with fixed inventories and signed/complex activities. Our cellwise
exponential class is not that covering class and is not proved invariant
under a gauge-theory RG step. Before claiming a cutoff comparison, one
must establish a suitable positive endpoint class, prove its boundary
marginal or energy-moment control with uniform constants, and treat the
exact covering identities, residual inventories and all eligibility queries.
Signed restricted sums must not be used as normalizers.

A concrete next deliverable is to test one explicitly specified
coarse-graining step or covering-weight construction against these two
requirements: a local edit cost and a uniform exceptional-label estimate
in its \emph{positive complete endpoint measure}. If elementary-cell
support or the declared reflection structure is lost, prove a replacement
hypothesis instead of assuming closure. Composition involving retained
exceptional labels and adaptive histories also requires a joint bound
rather than repeated marginal tails. The fixed-physical-size defect test,
nontrivial continuum fields and reconstruction remain downstream.

\section{Checks and scope of verification}
The accompanying standard-library script tests exact finite conditional
marginals, normalized reweighting, local likelihood norms, input/output
weight distinctions, domination and composition of accepted kernels,
cell incidence and reflection geometry, and noncommuting quaternion-group
link identities. A two-dimensional finite $\mathbb Z_2$ torus is enumerated
exactly for weighted event chessboard diagnostics. These finite models do
not calculate the four-dimensional $\SU$ probabilities and do not prove
universal analytic assertions. The earlier Python checks and standalone
Lean scalar certificates are retained and rerun; no new analytic or gauge
result is Lean-formalized. Independent review remains outstanding.

\begin{thebibliography}{9}
\bibitem{third}
\emph{Marginal boundary energy costs in the Wilson model}, third-pass
local research note, October 9, 2026,
\texttt{../boundary\_cost/boundary\_cost.tex}. Reviewed source SHA-256 is
recorded in \texttt{audit\_third\_pass.md}.
\bibitem{bk}
M.~Biskup and R.~Koteck\'y,
\emph{Forbidden gap argument for phase transitions proved by means of
chessboard estimates}, Commun. Math. Phys. \textbf{264} (2006), 631--656,
arXiv:math-ph/0505011, author manuscript \S2.2, Theorem 2.4, (2.10).
\bibitem{notes}
M.~Biskup, \emph{Reflection positivity and phase transitions in lattice
spin models}, author lecture notes, \S5.3, Theorem 5.8, (5.35), pp.~43--46.
\bibitem{o3}
OpenAI, \emph{The canonical massive continuum limit of the
two-dimensional $O(3)$ model}, October 4, 2026, supplied repository,
\texttt{build/sections/comparison.tex}. Inspected comparison requirements;
no transfer of its continuum result to the Wilson model is assumed.
\end{thebibliography}
\end{document}
